\section{Related Work and Literature Survey}
\label{sec:related}
\system{} sits at the intersection of ten research areas. We review each,
identify representative methods, their strengths and limitations, and how
the implemented architecture relates to them. Table~\ref{tab:comparison}
consolidates the comparison against representative campus-oriented and
RAG-oriented systems.

\subsection{Retrieval-Augmented Generation}
RAG~\cite{lewis2020rag} conditions a generative model on passages retrieved
from a non-parametric memory, improving factuality on knowledge-intensive
tasks and allowing the knowledge store to be updated without retraining.
The now-standard pipeline --- chunk a corpus, embed the chunks, retrieve
top-$k$ by vector similarity, and prompt an LLM with the concatenated
passages --- is surveyed comprehensively by Gao~\etal~\cite{gao2023ragsurvey},
who taxonomize \emph{naive}, \emph{advanced}, and \emph{modular} RAG and
catalogue the failure points: retrieval misses, low-precision context,
context-conflict, and over-generation beyond the evidence. Shuster
\etal~\cite{shuster2021retrieval} show that retrieval augmentation
measurably reduces hallucination in dialogue but does not eliminate it,
motivating additional grounding safeguards. Agentic variants such as
ReAct~\cite{yao2023react} interleave retrieval with reasoning steps.

\emph{Relation to \system{}.} The implemented pipeline is ``advanced RAG''
in Gao~\etal's taxonomy: it adds pre-retrieval query expansion, a
post-retrieval reranking and context-selection stage, a confidence gate,
and a post-generation verification stage around the core retrieve-then-read
loop. It is deliberately \emph{not} agentic in the ReAct sense --- the LLM
is never given tools or a multi-step scratchpad --- because the target
deployment is a low-resource, CPU-only local model where each generation
call costs several seconds and predictable single-shot latency was judged
more important than iterative refinement.

\subsection{Dense Retrieval and Sentence Embeddings}
Dense passage retrieval (DPR)~\cite{karpukhin2020dpr} trains a dual-encoder
to embed queries and passages into a shared space where relevance is inner
product, outperforming BM25 on open-domain QA when in-domain training data
is available. Sentence-BERT~\cite{reimers2019sbert} makes such embeddings
practical for unsupervised semantic similarity by fine-tuning
BERT~\cite{devlin2019bert} --- itself a Transformer
encoder~\cite{vaswani2017attention} --- with a siamese objective. Earlier
static embeddings~\cite{mikolov2013word2vec,pennington2014glove} lacked the
contextualization such tasks need. The widely used \code{all-MiniLM-L6-v2}
model combines the Sentence-BERT recipe with the MiniLM distillation of
Wang~\etal~\cite{wang2020minilm} to yield a 384-dimensional encoder small
enough to run on commodity CPUs.
Late-interaction models such as ColBERT~\cite{khattab2020colbert} retain
token-level detail at higher index cost. Approximate nearest-neighbor
indexes --- HNSW~\cite{malkov2020hnsw}, the graph structure used by the
ChromaDB store in this work, and IVF/PQ methods~\cite{johnson2019faiss} ---
make vector search scale.

\emph{Relation to \system{}.} Dense retrieval is used as an
\emph{adopted component}, not a contribution: the pretrained
\code{all-MiniLM-L6-v2} encoder is run locally with L2-normalized
embeddings and cosine ANN search over an HNSW-indexed ChromaDB collection.
No fine-tuning is performed, because no in-domain (query, passage) relevance
data exists for this institution.

\subsection{Lexical Retrieval and BM25}
The probabilistic relevance framework and its BM25 instantiation
\cite{robertson1994okapi,robertson2009bm25} remain the strongest
unsupervised sparse baseline. BM25 scores a document by a saturating
term-frequency function, an inverse-document-frequency weight derived from
Sp\"arck~Jones's term-specificity argument~\cite{sparckjones1972idf}, and a
length-normalization term controlled by parameters $k_1$ and $b$
\cite{manning2008ir}. Its strength is exact-term precision --- it never
confuses ``CSE'' with ``ECE'' the way a smooth embedding space can --- and
its weakness is vocabulary mismatch: a paraphrased query with no shared
terms retrieves nothing.

\emph{Relation to \system{}.} BM25 (the Okapi variant, \code{rank\_bm25})
is indexed over the same chunk set as the dense store, with a deliberately
minimal lowercase-alphanumeric tokenizer and no stemming or stop-word
removal, so that acronyms, department codes, and room numbers survive
intact. It is one of two retrieval signals, never used alone.

\subsection{Hybrid Retrieval and Reranking}
Combining dense and sparse signals consistently beats either alone
\cite{wang2021interpolation}. Fusion can be score-level (weighted sums of
normalized scores) or rank-level (reciprocal rank fusion,
RRF~\cite{cormack2009rrf}); Bruch~\etal~\cite{bruch2023fusion} analyze
convex combinations and show that normalization choice matters more than
the exact weight. A second stage typically reranks a shortlist with a
cross-encoder~\cite{nogueira2019bertrerank} or a sequence-to-sequence
model~\cite{nogueira2020monot5}, trading latency for precision; learning-
to-rank with support-vector regression~\cite{smola2004svr} is a lighter
alternative when labeled relevance judgments exist.

\emph{Relation to \system{}.} The implemented fusion is a convex score-level
combination ($0.6$ dense, $0.4$ lexical) over per-method min--max
normalized candidate sets, consistent with Bruch~\etal's finding that
per-method normalization is what makes a fixed weight meaningful. The
second stage is a \emph{deterministic four-feature weighted sum}, not a
neural cross-encoder, chosen for explainability and zero added model cost.
A support-vector reranker is implemented behind the same interface but is
\emph{never trained} for lack of relevance labels, and is reported as an
unevaluated extension point rather than a result.

\subsection{Hallucination Mitigation and Grounded Generation}
Ji~\etal~\cite{ji2023hallucination} survey intrinsic and extrinsic
hallucination and the mitigation families: faithful decoding, retrieval
grounding, and post-hoc verification. Post-hoc approaches include
SelfCheckGPT~\cite{manakul2023selfcheckgpt}, which samples multiple
generations and checks consistency, and FActScore~\cite{min2023factscore},
which decomposes a generation into atomic facts and verifies each against a
knowledge source. RAGAS~\cite{es2024ragas} operationalizes
\emph{faithfulness}, \emph{answer relevancy}, and \emph{context relevancy}
as LLM-judged metrics for RAG systems.

\emph{Relation to \system{}.} The implemented design layers four
independent safeguards: (i)~retrieval grounding; (ii)~a grounding-only
system prompt with explicit negative instructions; (iii)~a confidence gate
that returns a fixed refusal string \emph{without invoking the LLM} when
the retrieval-quality signal is low; and (iv)~a deterministic
post-generation pass that extracts phone-number, currency, room-number, and
year patterns from the generated answer, reduces each to a digit sequence,
and withholds the whole answer if any such sequence is absent from the
retrieved context. This last check is narrower than FActScore (it targets
checkable numeric claims, not arbitrary sentences) but requires no
additional model call --- an appropriate trade-off for a local-inference
deployment. Formal faithfulness measurement in the style of RAGAS is
proposed but not yet performed (Section~\ref{sec:eval}).

\subsection{Multi-Agent Systems and Query Routing}
The multi-agent paradigm~\cite{wooldridge2009mas} decomposes a task across
specialized agents coordinated by a supervisor or a shared protocol. LLM-
based instantiations --- AutoGen~\cite{wu2023autogen},
MetaGPT~\cite{hong2024metagpt}, HuggingGPT~\cite{shen2023hugginggpt}, and
generative agents~\cite{park2023generative} --- use an LLM both to route
and to execute, which is powerful but expensive and hard to audit. Intent-
based routing, by contrast, classifies the query once and dispatches to a
handler; the classifier can be rule-based, a shallow neural model, or an
LLM.

\emph{Relation to \system{}.} The router is a \emph{deterministic}
phrase- and keyword-matching supervisor: a $23$-entry navigation-phrase
list, a bare-room-number regular expression, and four domain keyword sets
are checked in order, and a trained LSTM intent classifier is consulted
\emph{only} when no deterministic rule fires. Every routing decision is a
literal substring match logged with its trigger, satisfying an explicit
auditability requirement. This is a design choice against the LLM-router
trend, motivated by predictability, zero routing latency, and the ability
to explain any decision post hoc. A lightweight ``multi-domain'' path
splits conjunctive queries into clauses and dispatches each independently.

\subsection{Intent Classification with Recurrent Models}
Spoken-language-understanding research established RNN-based joint intent
detection and slot filling on corpora such as ATIS~\cite{hemphill1990atis}
and SNIPS~\cite{coucke2018snips}; attention-augmented
bi-LSTMs~\cite{liu2016attention} and later BERT-based
classifiers~\cite{chen2019bertintent} are the standard baselines. LSTMs
\cite{hochreiter1997lstm} remain a reasonable choice for short-utterance
classification when training data is scarce and inference must be cheap,
though small datasets invite overfitting~\cite{goodfellow2016deep}. Open
instruction-tuned LLMs~\cite{touvron2023llama2,grattafiori2024llama3} are
an alternative zero-shot classifier but incur per-query inference cost.

\emph{Relation to \system{}.} The implemented classifier is an
\code{embedding}$\rightarrow$\code{LSTM}$\rightarrow$\code{dropout}$\rightarrow$\code{linear}
network with $32$-dimensional embeddings and hidden state, trained on
$168$ hand-authored synthetic utterances across $12$ intent classes. Its
measured held-out accuracy ($0.529$, against a training accuracy of
$1.00$) is characteristic of overfitting on a tiny synthetic set; we
therefore use it strictly as a \emph{fallback signal} behind the
deterministic rules and report its confusion matrix honestly
(Section~\ref{sec:results}).

\subsection{Virtual Tours and 360$^{\circ}$ Panoramic Navigation}
Immersive 360$^{\circ}$ tours are widely used in cultural heritage and
real estate~\cite{argyriou2020360design}, and hyperlinked panoramic video
navigation~\cite{neng2010hypervideo} predates them. The dominant
consumer reference is Google Street View~\cite{anguelov2010streetview},
which models the world as a graph of panoramas connected by navigable
links and preserves camera heading across transitions to maintain a sense
of continuous motion. Web viewers such as
Pannellum~\cite{petroff2019pannellum} render equirectangular imagery in the
browser.

\emph{Relation to \system{}.} The tour is a Street-View-style panorama
graph, but with an explicit data-structure contribution: each floor is a
\emph{doubly linked list} of scene nodes (enabling $O(1)$ sequential
traversal and ``scene $x$ of $y$'' counting scoped to a floor), branch and
vertical connections are held as separate cross-reference pointers, and a
dedicated construction pass links each floor's tail node to the next
floor's head node so that sequential traversal --- manual or automatic ---
crosses floor boundaries with no boundary-aware code. Entry heading is
carried on each link, Street-View style.

\subsection{Indoor Navigation, Wayfinding, and A$^{\star}$}
Indoor navigation systems~\cite{fallah2013indoor} combine a localization
technology~\cite{zafari2019localization} with a routing algorithm over a
navigation graph. The routing algorithm is almost always
Dijkstra's~\cite{dijkstra1959note} or its informed generalization
A$^{\star}$~\cite{hart1968astar}, which expands nodes in order of
$f(n)=g(n)+h(n)$ and is optimal when the heuristic $h$ is admissible
\cite{russell2020aima}.

\emph{Relation to \system{}.} The engine implements A$^{\star}$ with a
binary heap over an adjacency-list graph built on demand from the
relational schema. Its heuristic degrades in three tiers --- Euclidean
distance on planar floor-plan coordinates, an equirectangular great-circle
approximation on latitude/longitude, or zero --- and at the zero tier the
search is provably equivalent to Dijkstra's, which is used directly for
single-source nearest-facility queries. Localization is out of scope: the
endpoint consumes graph node identifiers, and mapping a physical position
to a node is left to the client.

\subsection{Conversational AI for Educational Institutions}
Systematic reviews of educational chatbots
\cite{okonkwo2021chatbots,perez2020chatbots,winkler2018chatbots} report
that most deployed systems are intent-and-response or FAQ retrieval
systems, that admissions and student-services queries dominate real usage,
and that trust and answer correctness are the principal barriers to
adoption. Few reported systems combine a conversational assistant with
spatial exploration of the same campus.

\emph{Relation to \system{}.} \system{} targets exactly the
admissions/academics/facilities query mix these reviews identify, but
replaces the FAQ-retrieval or scripted-intent backend with grounded RAG
over the institution's own documents, and couples it to a spatial tour and
a wayfinding graph of the same campus --- the integration that the reviews
find largely absent.

\begin{table*}[!t]
\centering
\caption{Positioning of \system{} Against Representative Campus-Oriented and RAG-Oriented Systems.
Entries describe \emph{category} characteristics; \system{} entries describe the \emph{implemented} artifact.}
\label{tab:comparison}
\footnotesize
\renewcommand{\arraystretch}{1.25}
\begin{tabular}{@{}p{2.7cm}p{2.1cm}p{2.0cm}p{1.9cm}p{1.6cm}p{1.7cm}p{3.0cm}@{}}
\toprule
\textbf{System class} & \textbf{Info.\ retrieval} & \textbf{Generation / QA} & \textbf{Multi-agent routing} & \textbf{360$^{\circ}$ tour} & \textbf{Graph wayfinding} & \textbf{Principal limitation} \\
\midrule
Institutional website + scripted chatbot &
Keyword / FAQ match &
Templated / canned &
No &
No &
No &
No paraphrase robustness; no spatial context \\

Educational FAQ chatbot \cite{okonkwo2021chatbots,perez2020chatbots} &
Intent + slot &
Retrieval of stored answers &
Intent classifier only &
No &
No &
Answer set is fixed; brittle to unseen intents \\

Naive RAG assistant \cite{lewis2020rag,gao2023ragsurvey} &
Dense top-$k$ &
LLM over concatenated passages &
No &
No &
No &
Hallucinates on weak/thin retrieval; no refusal \\

Agentic RAG (ReAct-style) \cite{yao2023react,wu2023autogen} &
Iterative dense &
LLM with tool loop &
LLM-as-router &
No &
No &
High latency; routing not auditable; cost \\

360$^{\circ}$ virtual tour product \cite{argyriou2020360design,petroff2019pannellum} &
None &
None &
No &
Yes (scenes + hotspots) &
No &
No queryable knowledge; no routes \\

Street View-style panorama graph \cite{anguelov2010streetview} &
None &
None &
No &
Yes (panorama graph) &
Partial (link graph) &
Outdoor scale; no institutional knowledge \\

Indoor navigation system \cite{fallah2013indoor,zafari2019localization} &
None &
None &
No &
No &
Yes (Dijkstra / A$^{\star}$) &
Requires a localization technology; no QA \\

\midrule
\textbf{\system{} (this work)} &
\textbf{Hybrid dense + BM25, $0.6/0.4$ fusion, 4-feature rerank} &
\textbf{Local Llama~3.2, grounding-only prompt, confidence gate + post-hoc numeric check} &
\textbf{Deterministic phrase/keyword first; LSTM fallback; logged} &
\textbf{Yes: per-floor doubly linked list, cross-floor stitching, guided walk} &
\textbf{Yes: A$^{\star}$ over live relational graph, 3-tier heuristic} &
\textbf{No formal evaluation yet; A$^{\star}$ has no front-end; SVR reranker untrained} \\
\bottomrule
\end{tabular}
\end{table*}

